
% !TEX TS-program = pdflatex
\documentclass[11pt]{article}

\usepackage[T1]{fontenc}
\usepackage[utf8]{inputenc}
\usepackage{lmodern}
\usepackage{microtype}
\usepackage{geometry}
\geometry{margin=1in}

\usepackage{amsmath,amssymb,amsthm}
\usepackage{mathtools}
\usepackage{stmaryrd}
\usepackage{hyperref}
\usepackage{enumitem}
\usepackage{booktabs}

\hypersetup{
  colorlinks=true,
  linkcolor=blue,
  urlcolor=blue,
  citecolor=blue
}

% --- Simple macros ---
\newcommand{\OSLF}{\textsf{OSLF}}
\newcommand{\NTT}{\textsf{NTT}}
\newcommand{\NT}{\textsf{NT}}
\newcommand{\MeTTaIL}{\textsf{MeTTaIL}}
\newcommand{\rhocalc}{\textsf{\(\rho\)}-calculus}
\newcommand{\picalc}{\textsf{\(\pi\)}-calculus}
\newcommand{\GF}{\textsf{GF}}
\newcommand{\WM}{\textsf{WM}}

\newcommand{\BB}{\mathbb{B}}
\newcommand{\Set}{\mathbf{Set}}
\newcommand{\C}{\mathcal{C}}
\newcommand{\D}{\mathcal{D}}

\newcommand{\bx}{\Box}
\newcommand{\dia}{\Diamond}

\newcommand{\leanfile}[1]{\texttt{#1}}
\newcommand{\leanid}[1]{\texttt{#1}}

\title{Algorithmic Spatial--Behavioral Types from Rewrite Systems\\
\large A Primer on \OSLF{}, Native Type Theory, and Evidence-Valued Extensions\\
\small (Conceptual notes intended to accompany a Lean formalization snapshot; 2026-02-25)}
\author{}
\date{}

\begin{document}
\maketitle

\begin{abstract}
\OSLF{} (Operational Semantics in Logical Form) proposes an algorithm that takes a language
presented as a rewrite system---syntax with binders, structural equations, and reduction rules---and
\emph{mechanically generates} a spatial--behavioral logic and an associated type/judgment system.
The construction is category-theoretic: starting from a category encoding syntax and dynamics,
one passes through a presheaf topos via Yoneda, recovers internal Heyting structure from subobjects,
and applies a Grothendieck construction to obtain a category of \emph{native types}.
\NTT{} (Native Type Theory) describes a corresponding internal language with dependent types,
logical connectives, and (in principle) inductive/coinductive fixed points.

This note distills the conceptual picture with an eye toward formalization and executability.
In addition to the classical ``types as predicates'' reading, we highlight a \emph{valuation layer}:
OSLF propositions can be interpreted not only as truth values but also as evidence values
(e.g. in a Heyting/quantale/fuzzy setting), enabling probabilistic, paraconsistent, or resource-aware
extensions without changing the underlying operational core.

The document separates: (i) claims that are standard in the literature, (ii) endpoints that are already
formalized in the accompanying Lean snapshot, and (iii) plausible research directions that are not yet
theorems.
\paragraph{Status note.} Implementation-status and count claims in this draft are snapshot-scoped; verify against current repository trackers and build outputs before citing as current state. Lean-proven claims are only those tied to explicit Lean file/theorem references; broader mathematical or historical claims are literature-backed context and not asserted as Lean-proven unless explicitly stated.
\end{abstract}

\tableofcontents

\section{Why ``logic from semantics''?}

The usual workflow in programming languages is: define a calculus, then separately design (or port)
a logic and type system to reason about it. \OSLF{} reverses the order of invention: given a
machine-readable description of a model of computation as a rewrite system, it aims to
\emph{derive} a spatial--behavioral logic and type discipline by a uniform recipe.

Two guiding intuitions:

\begin{enumerate}[label=(\arabic*)]
\item \textbf{Types as queries over behavior.}
A type is not merely a shape constraint on syntax; it is a property that can talk about
structure (\emph{spatial} connectives) and dynamics (\emph{behavioral/modal} connectives).

\item \textbf{Witnesses matter.}
Under Curry--Howard and realizability-style readings, semantics of formulae are collections of
witnesses. Sets are one option, but not the only option; changing the witness-collection algebra
changes the induced logical connectives.
\end{enumerate}

For concurrency, the intended payoff is a uniform method to generate logics that can express
liveness/safety properties and structural facts (e.g.\ parallel composition) that are invisible to
many traditional modal logics.

\section{Input data: rewrite systems as ``graph-structured theories''}

\subsection{Rewrite-system presentation}
An \OSLF{} input language can be thought of as:

\begin{itemize}
\item \textbf{Syntax with binders:} term constructors, potentially higher-order, described using
a higher-order abstract syntax (HOAS) discipline.
\item \textbf{Equations:} a structural congruence/definitional equality that quotients raw terms.
\item \textbf{Reduction rules:} a (possibly labeled) transition relation generated from rules, often
context-closed and compatible with equations.
\end{itemize}

In the \rhocalc{} and related calculi, the mutual recursion between \emph{names} and \emph{processes},
and the associated quote/drop structure, is a canonical example: names can be built from processes,
and processes can execute names.

\subsection{From reductions to graph structure}
A recurring theme in \OSLF{} and related work is to represent operational semantics internally
as a \emph{graph structure} whose edges witness allowed transitions.
This is closely aligned with the view that operational semantics correspond to GSOS-style rules and
can be seen categorically as distributive laws.

\section{The \OSLF{} construction in one diagram}

\subsection{The conceptual pipeline}
A compact high-level view is:

\[
\text{(rewrite system)} \;\leadsto\;
\underbrace{\text{(category \(\C\) of syntax/constructors)}}_{\text{``constructor category''}}
\;\xrightarrow{\text{Yoneda}}\;
\widehat{\C} = \Set^{\C^{op}}
\;\leadsto\;
\text{(subobjects \(\Rightarrow\) Heyting structure)}
\;\xrightarrow{\text{Grothendieck}}\;
\text{native types } \NT(\C).
\]

In prose:

\begin{enumerate}[label=(\arabic*)]
\item Start from a category \(\C\) that encodes the language's sorts and (crucially) the
\emph{sort-crossing constructors} that move between sorts (e.g.\ quote/drop).
\item Apply Yoneda to move into a presheaf setting, where each object has a representable presheaf.
\item Interpret predicates as subobjects (subfunctors) of representables; the poset of subobjects
carries Heyting-algebra structure (internal logic).
\item Apply the Grothendieck construction to obtain a category whose objects can be read as
pairs \((X,U)\): a \emph{sort} \(X\) together with a \emph{predicate/filter} \(U\) on its representable,
i.e.\ a native type.
\end{enumerate}

\subsection{What is a native type?}
A useful operational reading is:

\begin{quote}
A \emph{native type} is a \emph{sort} together with a \emph{behavioral predicate} over terms/patterns of that sort.
\end{quote}

In the Lean snapshot, this ``type = (sort, predicate)'' idea appears explicitly in the generated-typing layer.

\subsection{Modalities and a Galois connection}
A distinctive feature of \OSLF{} is that modal operators are not bolted on: they are \emph{generated}.
In a two-sort calculus like the \rhocalc{}, sort-crossing constructors provide canonical candidates
for modal operators. Abstractly, one gets an adjunction (often presented as a Galois connection)
\(\dia \dashv \bx\) linking possibility and necessity.

Conceptually:

\begin{itemize}
\item \(\dia \varphi\) expresses ``there exists a behavior step (or path) after which \(\varphi\) holds''.
\item \(\bx \varphi\) expresses ``for all behavior steps (or predecessors, depending on polarity),
\(\varphi\) holds''.
\end{itemize}

The key point is structural: the adjunction is derived from the same categorical data used to
describe syntax and operational edges.

\section{Native Type Theory as internal language}

\subsection{Internal logic and dependent types}
\NTT{} describes the internal language of the native-type construction, including:

\begin{itemize}
\item propositional connectives (interpreted in a Heyting setting);
\item dependent sums \(\Sigma\) and products \(\Pi\), subsuming existential and universal quantification;
\item comprehension types \(\{x:A \mid \varphi\}\) connecting predicates and types.
\end{itemize}

The conceptual advantage is that once the presheaf/topos infrastructure is in place,
many logical constructions come from general categorical principles rather than ad hoc encodings.

\subsection{Fixed points: \(\mu/\nu\) as inductive/coinductive types}
A longer-term goal is to connect operational modal reasoning to a \(\mu\)-calculus style
fixed-point logic:

\begin{itemize}
\item \(\mu\) for least fixed points (inductive properties like reachability);
\item \(\nu\) for greatest fixed points (coinductive properties like invariants/bisimulation).
\end{itemize}

This is where Knaster--Tarski style fixed-point constructions in an appropriate fibration
meet the operational layer.

\section{From truth to evidence: valuation layers}

\subsection{Why evidence-valued semantics?}
In many applications (world modeling, uncertain inference, inconsistent data), one wants a semantics
more refined than \(\{\texttt{true},\texttt{false}\}\). One can keep the \OSLF{} operational core (the transition system)
and change only the \emph{codomain} of formula evaluation:

\[
\llbracket \varphi \rrbracket(p) \in V
\]
where \(V\) is a partially ordered algebra of evidence/weights rather than \(\BB\).

This makes the ``collection of witnesses'' parameter explicit.

\subsection{Residuated/tensor structure}
A common pattern is to endow \(V\) with a monoidal tensor \(\otimes\) and a residuation
(implication-like operator) \(\Rightarrow\), yielding resource-aware entailment:

\[
a \otimes b \le c \quad\Longleftrightarrow\quad b \le (a \Rightarrow c).
\]

In a fuzzy/probabilistic instance, one can take \(V = [0,1]\) with a chosen t-norm as \(\otimes\).
In a paraconsistent instance, \(V\) may track positive/negative support separately.

\subsection{Quantifiers as infimum/supremum}
When quantifiers are added, a clean semantics in a complete lattice is:

\[
\llbracket \forall x.\,\varphi \rrbracket = \bigwedge_{d \in \mathrm{Dom}} \llbracket \varphi[x:=d] \rrbracket,
\qquad
\llbracket \exists x.\,\varphi \rrbracket = \bigvee_{d \in \mathrm{Dom}} \llbracket \varphi[x:=d] \rrbracket.
\]

This mirrors the \(\Pi/\Sigma\) view in \NTT{}, and it is compatible with evidence lattices
beyond Booleans.

\section{Process calculi: \(\pi\), \(\rho\), and correspondence}

\subsection{Why correspondence matters}
Completeness and expressiveness results for modal/spatial logics are often proved
\emph{relative to} an operational notion of equivalence, such as (weak) bisimulation or barbed congruence.
To transfer results between calculi (e.g.\ from \picalc{} to \rhocalc{}), one typically needs a high-quality encoding
and an operational correspondence theorem: forward simulation and some form of backward reflection.

\subsection{Derived vs core semantics}
A recurring engineering issue in formalizations is separating:

\begin{itemize}
\item \textbf{core/canonical reduction:} the intended semantics of the calculus;
\item \textbf{derived/administrative steps:} convenience rules (e.g.\ unfolding replication, name admin)
that should not change observable behavior.
\end{itemize}

A robust architecture proves correspondence in the derived layer (because proofs are shorter),
then transports results back to the core via conservativity/adequacy theorems on a well-defined fragment.

\subsection{Substitutability as a ``crown'' statement}
A central claim in the \OSLF{} narrative is a substitutability theorem:
operational equivalence (bisimilarity) coincides with logical equivalence in the generated type logic.
Such theorems are often in the spirit of Hennessy--Milner results, but strengthened by spatial operators
that can observe structure.

\section{Linguistic and neural connections (speculative but testable)}

This note also records a line of thought motivated by grammar formalisms and neural language models.

\subsection{Operational cores of grammar}
One can view a grammar formalism as a rewrite system: it generates structures and admits reductions/derivations.
If multiple natural-language grammars are treated as instances, one can ask for a \emph{weakest operational core}
that preserves a shared logical interface (what can be expressed/queried).

This echoes ``minimalism'' intuitions in linguistics, but the formal content is:
\emph{identify a minimal rewrite kernel that conservatively preserves the induced spatial--behavioral logic.}

\subsection{Tensorization and distributional semantics}
A separate (and still open) research direction is to compile logical/operational forms into tensorial expressions
(using a tensor logic or related categorical semantics), thereby connecting symbolic derivations to
distributional/statistical models.

A pragmatic hypothesis for transformer-style architectures is that long-term memory need not store complete
sequence detail; in many tasks it suffices to store a structured partial order or summary that supports the
queries of interest. This is naturally phrased using order/enrichment (lattices/quantales) rather than raw lists.

\section{Lean formalization snapshot: what is already there}

The accompanying Lean snapshot (in a \texttt{mettapedia} project) contains a significant amount of
executable and theorem-proved infrastructure. This section only lists \emph{representative} anchor points,
without attempting to be complete.

\subsection{Executable reduction and soundness scaffolding}
A key milestone for usability is having an executable one-step reducer for a core calculus
with a proven soundness theorem, plus regression tests.

\subsection{Language-parametric generation}
The formalization factors out a \leanid{LanguageDef} describing a user-defined language
(constructors + rewrite rules), then mechanically builds:

\begin{itemize}
\item a generated native-type notion (\leanid{GenNativeType});
\item a generated typing judgment (\leanid{GenHasType});
\item a substitutability-style stability story (types preserved under substitution).
\end{itemize}

\subsection{Categorical bridge endpoints}
The snapshot also includes explicit endpoints that tie the categorical constructions to
concrete modal operators and internal logic packages.

\subsection{Evidence/tensor endpoints}
Finally, there are small bridge theorems for evidence/tensor structure (e.g.\ residuation on \([0,1]\))
intended to make it easy to connect the logic layer to quantitative interpretations.

\section{Claim hygiene and open problems}

\subsection{What is standard}
The following are standard (non-controversial) mathematical facts used by the program:

\begin{itemize}
\item Presheaf categories are toposes and support an internal Heyting logic.
\item Subobjects of a fixed object form a Heyting algebra (and often a complete lattice under mild conditions).
\item Grothendieck constructions package indexed families (fibrations) into total categories.
\end{itemize}

\subsection{What is formalized (snapshot-dependent)}
The Lean snapshot provides concrete instantiations of:
generated modal operators, Galois connections, internal-language bridge packages,
and executable reduction/matching infrastructure.
(See the module pointers in the appendix.)

\subsection{What remains genuinely open}
Some ambitious directions are not ``for free'' and should be treated as research problems:

\begin{itemize}
\item \textbf{Generic completeness of spatial logics.}
Even if modalities are generated functorially, completeness (in the sense of
``logical equivalence implies bisimulation'') typically demands image-finiteness or similar hypotheses.
\item \textbf{Transport of completeness across encodings.}
Transferring completeness from one calculus to another requires encodings that preserve
the relevant observation structure and adequate reflection properties.
\item \textbf{Fixed points + operational distributive laws.}
It is plausible that distributive-law structure helps organize proofs, but a true
``completeness-for-free'' theorem would need very careful statement and nontrivial conditions.
\end{itemize}

\appendix

\section{Appendix: suggested code-map anchors}

The following filenames/theorem identifiers are useful anchors when navigating the Lean snapshot:

\begin{itemize}
\item \leanfile{Mettapedia/OSLF/Framework/GeneratedTyping.lean} --- generated typing pipeline.
\item \leanfile{Mettapedia/OSLF/Framework/CategoryBridge.lean} --- modal adjunction endpoints.
\item \leanfile{Mettapedia/OSLF/Framework/ToposTOGLBridge.lean} --- topos/internal-language bridge packages.
\item \leanfile{Mettapedia/OSLF/QuantifiedFormula.lean} --- quantified formula extension and inf/sup semantics.
\item \leanfile{Mettapedia/Languages/ProcessCalculi/*} --- \(\rho\)/\(\pi\) calculi, encodings, bisimulation layers.
\end{itemize}

\section{Appendix: a minimal glossary}

\begin{description}[leftmargin=!,labelwidth=3.0cm]
\item[Rewrite system] Syntax + equations + reduction rules.
\item[Spatial logic] Logic with connectives that observe structure (e.g.\ parallel composition).
\item[Behavioral logic] Modal logic observing operational steps (e.g.\ possibility/necessity).
\item[Presheaf topos] Category of contravariant set-valued functors \(\Set^{\C^{op}}\).
\item[Heyting algebra] Lattice with implication appropriate for intuitionistic logic.
\item[Quantale] Complete lattice with monoidal product distributing over joins (resource/evidence models).
\item[Galois connection] Adjunction between posets; here used for \(\dia \dashv \bx\).
\end{description}

\vspace{1em}
\noindent\textbf{Acknowledgment of sources:} This primer is a synthesis of standard category-theoretic semantics,
the \OSLF{} and \NTT{} papers, and a Lean formalization snapshot in the \texttt{mettapedia} codebase.

\end{document}
